\documentclass[12pt]{article}
\usepackage[T1]{fontenc}
\usepackage[margin=1in]{geometry}
\usepackage{amsmath,amssymb,booktabs,array,enumitem}
\usepackage{graphicx,float}
\usepackage{textcomp}
\usepackage[colorlinks=true,linkcolor=blue,citecolor=blue]{hyperref}
\def\bxthree{BX3}
\def\purpose{\textit{Purpose Layer}}
\def\bounds{\textit{Bounds Engine}}
\def\fact{\textit{Fact Layer}}

\begin{document}

\section{Background}

The proliferation of autonomous multi-agent AI systems has exposed a fundamental tension between execution speed and accountability. As autonomous agents assume operational roles previously reserved for human actors—scheduling, resource allocation, environmental intervention—the question of who bears responsibility when outcomes diverge from intent has become architecturally urgent. This section situates the Bailout Protocol within four intersecting research domains: accountability in multi-agent AI systems, human-in-the-loop (HITL) design, fail-safe design in safety-critical systems, and the BX3 Framework's three-layer model for functional role decomposition.

\subsection{Accountability in Multi-Agent AI Systems}

The problem of accountability in distributed autonomous systems is not new, but its stakes have grown with the operational scope of AI agents. Early multi-agent systems~\cite{trist1981} were designed around modular task decomposition: a parent agent spawns child agents to handle subtasks, with communication protocols governing inter-agent messaging. However, this architectural model inherits a fundamental weakness: when a chain of delegated agents produces an unintended outcome, the causal chain between the root human principal and the terminal action becomes opaque.

This problem is compounded when agents operate with partial observability, asynchronous execution, and heterogeneous reasoning models. Bansal et al.\ \cite{bansal2019} identified a critical failure mode in sequential reasoning systems where early decisions constrain downstream possibilities in ways that are neither visible nor reversible to the originating authority. Their work on safe policy synthesis in sequential decision processes demonstrates that uncertainty compounding through a chain of reasoning steps produces emergent failure modes that no single agent in the chain can fully anticipate. When multiple agents contribute to a composite outcome—each acting within its local information state—the resulting global trajectory can be neither traced to a single decision point nor attributed to any individual human principal.

In conventional software systems, accountability is maintained through logging, auditing, and access control: mechanisms that presuppose a single authoritative execution environment. Multi-agent AI systems break this assumption because reasoning and execution are distributed across actor boundaries. The result is what has been termed the \textit{black box problem}: autonomous actions that are technically attributable to an organization but operationally orphaned from any individual human decision-maker. This failure mode is not merely an artifact of insufficient logging; it is structurally embedded in systems where reasoning has been distributed without a corresponding architectural commitment to preserving human accountability anchors at each delegation boundary. The Bailout Protocol addresses this failure mode not by improving logging granularity after the fact, but by ensuring that any unresolvable condition in the execution chain propagates upward to a designated human authority before any irreversible action is taken. The protocol is thus not a logging mechanism but a structural constraint on the delegation graph itself.

\subsection{Human-in-the-Loop Design}

Human-in-the-loop (HITL) architectures have long been recognized as essential for high-stakes AI deployments. The foundational insight, traceable to Wiener \cite{wiener1948} and his analysis of feedback-driven automatic systems, is that machines operating in uncertain environments must have their outputs subject to human review at critical decision boundaries. Without this, the system optimizes against a model that is necessarily incomplete relative to the operational reality it navigates. Wiener's cybernetic framework established that every effective automatic control loop must contain a human supervisory component—not as an optional failsafe but as a structurally necessary element when systems operate in environments characterized by irreducible uncertainty and adversarial perturbation.

Modern HITL frameworks \cite{bansal2019} decompose human oversight into three temporal positions relative to agent action: \textit{pre-loop} oversight (humans set constraints before execution), \textit{intra-loop} oversight (humans intervene during execution), and \textit{post-loop} oversight (humans audit outcomes after execution). Each position provides different guarantees: pre-loop oversight shapes the space of permissible actions; intra-loop oversight enables real-time course correction; post-loop oversight supports accountability and learning but cannot prevent harm that has already occurred. The architectural placement of the human within this temporal structure determines the strength of the accountability guarantee. Pre-loop oversight alone leaves the system fully autonomous between constraint-setting events; post-loop oversight alone cannot prevent irreversible harm. Only intra-loop oversight, where the human is structurally present during execution, provides a genuine hard constraint on the system's action space.

The BX3 Framework integrates all three temporal positions structurally through its three-layer model. The \purpose\ layer (Layer 1) encodes pre-loop oversight as immutable strategic constraints. The \bounds\ engine (Layer 2) operates within those constraints and reports to the \purpose\ layer, enabling intra-loop visibility. The \fact\ layer (Layer 3) generates the \textit{Outcome Record} (Plane 7 of the 9-Plane DAP architecture), which is the substrate for post-loop accountability. The Bailout Protocol extends this model by guaranteeing that any condition falling outside the \bounds\ engine's resolution capacity triggers an exception path that bypasses all Machine Accountability Anchors and surfaces directly to the Human Root, ensuring that the human decision-maker is always operationally present at the precise point where the system's model of the world fails.

The architectural distinction between the BX3 HITL model and conventional HITL implementations is critical to understand. Conventional HITL systems implement human review as a checkpoint or approval gate: an external constraint on the agent's execution path, analogous to a gate that must receive a human signal before the system proceeds. The BX3 model instead bakes human oversight into the functional architecture of every node. The \purpose\ layer is not a checkpoint; it is a structural component of each BX3 Loop, guaranteeing that the human anchor is always one layer below any Machine actor in the resolution chain, regardless of how far recursive spawning has propagated the system into edge environments. This means that HITL is not a mode that can be disabled or bypassed—it is a property of the layer structure itself.

\subsection{Fail-Safe Design in Safety-Critical Systems}

The design of fail-safe mechanisms in safety-critical systems has a long pedigree in computer science. Dijkstra's foundational work on structured programming and program correctness \cite{dijkstra1982} established the principle that the correctness of a system cannot be established by testing alone: correctness must follow architecturally from the structure of the computation. Applied to fail-safe design, this principle demands that systems fail in a predetermined, architecturally enforced manner rather than through undefined or emergent behavior. A system that fails arbitrarily—entering an undefined state from which its subsequent behavior is unpredictable—is not a fail-safe system; it is merely a system that has failed. The fail-safe property requires that the system's behavior upon encountering any unresolvable condition is specified in advance and enforced by the architecture, not by runtime estimation of what corrective action might seem appropriate.

In safety-critical control systems, the dominant fail-safe pattern is \textit{redundancy with majority voting} or \textit{degraded mode operation}, where multiple independent subsystems compute the same output and the system defaults to a safe state when consensus is not achieved. N-version programming, a related technique, executes independent implementations of the same specification and compares results to detect implementation-specific failures. These approaches improve reliability against random hardware failures and implementation bugs but do not address the failure mode that arises when the system encounters conditions outside its designed operational envelope—the \textit{unknown unknown} failure regime. In this regime, the system's own model of its operational boundaries is incomplete, which means the system's built-in failure responses are themselves parameterized for conditions that do not match the actual failure.

The Bailout Protocol operates in this distinct failure regime: it does not address random failure or implementation error, but rather the condition where a node exhausts its modeled resolution capacity. The standard response in safety-critical design is to fall back to a safe state—hold position, shut down, request input—but in distributed autonomous systems, this fallback requires a communication channel to a responsible authority. The Bailout Protocol's contribution is to guarantee that channel structurally: the protocol is embedded in every BX3 Loop's Worksheet, meaning that even during cloud disconnection (when communication to centralized systems is impossible), each node retains the complete Bailout Protocol and can escalate to the Human Accountability Anchor via the next available communication window. The architectural embedding of fail-safe logic at every node, rather than at centralized checkpoints, reflects the same design philosophy that governs safety-critical distributed control systems: the system must remain safe even when its communication infrastructure is partially degraded. The Bailout Protocol treats network connectivity not as a prerequisite for safety but as a resource whose availability varies, and designs the fail-safe mechanism to be robust across the full range of connectivity states.

\subsection{The BX3 Framework's Three-Layer Model}

The BX3 Framework \cite{bx3framework2026} proposes a functional decomposition of any autonomous node into three architecturally distinct layers, each with a specific role in maintaining system-wide deterministic integrity. This decomposition is not a software pattern but a formal architectural constraint: each layer has a distinct type of occupant, a distinct set of permissions, and a distinct relationship to the accountability chain.

\begin{itemize}
    \item \textbf{Layer 1 — Purpose (Human Accountability Anchor):} Sets strategic goals, authorization boundaries, and success criteria. The Purpose layer is the only layer that may be occupied by a Human actor; it is architecturally prohibited from residing inside a Machine-only execution context. The Purpose layer generates Plane 1 (Mandate) and Plane 2 (Intent) of the 9-Plane Decision Accountability Plane (DAP) architecture.
    \item \textbf{Layer 2 — Bounds Engine (Limbless Proposer):} Computes, proposes, and simulates state transitions within the boundaries set by the Purpose layer. The Bounds Engine has no direct execution permissions; it operates exclusively in the reasoning plane, generating Plane 4 (Reason Log), Plane 5 (Decision), and Plane 6 (Projection). Its outputs feed the Sandbox Gate.
    \item \textbf{Layer 3 — Fact Layer (Physical Firewall):} The only layer with execution authority. The Fact Layer applies hard-coded physical, regulatory, and safety constraints to any incoming command from the Bounds Engine and either executes or blocks based on constraint evaluation. It generates Plane 7 (Outcome Record) and Plane 8 (Execution), and validates Plane 9 (Projection Confirmation) through sandbox simulation.
\end{itemize}

The critical property of the three-layer model is \textit{interchangeability}: the functional properties (Purpose, Bounds Engine, Fact) are decoupled from specific actor implementations. A Human actor and a Machine actor can occupy the same functional layer, provided the layer's constraints are satisfied. This means the BX3 Framework can govern both human-operated and AI-operated nodes through the same architectural rules, enabling uniform accountability guarantees across heterogeneous execution contexts.

The Bailout Protocol is the exception mechanism of the BX3 Framework. When the Bounds Engine encounters a condition it cannot resolve within its modeled capacity—either because the condition falls outside its operational envelope or because required information is unavailable—the system does not attempt to estimate a corrective action. Instead, it generates a \textit{Cascading Trigger}: an exception event that carries the complete problem context upward through the recursive tree, bypassing all intermediate Machine actors, until it reaches the Human Accountability Anchor. This behavior is not implemented as a fallback handler; it is a structural property of the BX3 Loop architecture, embedded in the Worksheet delivered to every node during recursive spawning. The three-layer model also provides the theoretical foundation for the Bailout Protocol's exclusion of Machine Accountability Anchors from the resolution path: because the Purpose layer is architecturally required to remain Human-anchored, any Machine actor in the resolution chain is necessarily operating in a subordinate capacity. The protocol's bypass mechanism is therefore not an arbitrary design choice but a direct consequence of the layer separation: exceptions propagate to the highest-purpose authority in the chain, which by construction must be a Human actor.

\subsection{Summary}

The research landscape surrounding accountability in autonomous systems spans distributed AI architecture, HITL design principles, safety-critical systems engineering, and formal functional decomposition. Each domain contributes a necessary constraint: multi-agent systems must track causal chains through delegation layers; HITL systems must position human oversight at points of maximum leverage; fail-safe systems must define behavior at the boundaries of their operational envelope; and functional decomposition frameworks must maintain clean separation between reasoning and execution.

The BX3 Framework's three-layer model and the Bailout Protocol together constitute an architecture that satisfies all four constraints simultaneously. The three-layer model provides functional decomposition that makes reasoning/execution separation architecturally enforceable. The 9-Plane DAP provides the accountability substrate for post-hoc auditing. The Recursive Spawning mechanism with Worksheet embedding ensures that every node in the system carries the complete Bailout Protocol, making fail-safe behavior a property of each node rather than a centralized checkpoint. And the Cascading Trigger mechanism guarantees that unresolvable conditions always escalate to a Human Accountability Anchor, eliminating the black box failure mode where autonomous actions become permanently orphaned from human oversight.

The following sections of this paper elaborate the Bailout Protocol's formal specification, trigger conditions, escalation mechanics, and implementation within the BX3 Framework.

\end{document}
